\appendix

\section{Statement--Lean Parity Tables}\label{app:selected-crosswalk}\label{app:formal-boundaries}

\subsection{Current Lean correspondence}

\Cref{tab:paper1-native} maps the current coupled construction and recovery
results to their canonical Lean declarations in this repository. The table
includes the encryption and byte decoders, finite-field ciphertext
refinement, key reuse, and one-use random-oracle privacy. Write \lean{HE}
for \lean{HomotopicalEncryption}. The historical Palomar comparison
appears in \cref{app:historical-palomar}.

To build the native modules and reproduce both byte round trips, run the
following from \lean{lean/} with this repository's toolchain and dependencies.
\begin{verbatim}
lake build Protocol.Transfer Protocol.Storage.CoupledRecovery.Demo
lake env lean --run Protocol/Storage/CoupledRecovery/Demo.lean
\end{verbatim}

\begingroup
\small
\setlength{\tabcolsep}{6pt}
\renewcommand{\arraystretch}{1.12}
\begin{longtable}{@{}>{\raggedright\arraybackslash}p{0.39\linewidth}>{\raggedright\arraybackslash}p{0.57\linewidth}@{}}
\caption{HCP encryption construction and quantitative recovery}\label{tab:paper1-native}\\
\toprule
\textbf{Statement} & \textbf{Canonical Lean declarations}\\
\midrule
\endfirsthead
\caption[]{HCP encryption construction and quantitative recovery (continued).}\\
\toprule
\textbf{Statement} & \textbf{Canonical Lean declarations}\\
\midrule
\endhead
\bottomrule
\endfoot
\textbf{P1-N01. Intrinsic lattice.}\newline
\Cref{prop:coupled-harmonic-graph,prop:coupled-intrinsic-lattice}.
The integer chart covers the entire original intrinsic lattice, with basis,
rank, discreteness and full lattice span inside harmonic space.
&
\lean{HE.HCP.CoupledCVP.harmonicLatticePresentation}\par
\lean{HE.HCP.CoupledLattice.integerGraphEquiv_eq_point}\par
\lean{HE.HCP.CoupledLattice.intrinsic_finrank}\par
\lean{HE.HCP.CoupledLattice.intrinsic_discrete}\par
\lean{HE.HCP.CoupledLattice.intrinsicInHarmonicSpace_isZLattice}
\\\midrule
\textbf{P1-N02. Public class and operative geometry.}\newline
\Cref{thm:coupled-public-class} and the derived Hodge operator.
&
\lean{HE.HCP.CoupledSemantics.semanticEquivalence}\par
\lean{HE.HCP.CoupledGeometry.degreeOneLaplacian_coordinates}\par
\lean{HE.HCP.CoupledGeometry.coupling_eq_of_laplacianMatrix_eq}
\\\midrule
\textbf{P1-N03. Ciphertext and recovery diagrams.}\newline
\Cref{prop:coupled-native-transport,prop:coupled-recovery-square}.
&
\lean{HE.HCP.CoupledTransport.transport_observed}\par
\lean{HE.HCP.CoupledGeometry.harmonicOneProjection_transport}\par
\lean{HE.HCP.CoupledGeometry.information_projection_transport}
\\\midrule
\textbf{P1-N04. Exact finite recovery.}\newline
\Cref{prop:coupled-byte-recovery}; all-input decoder refinement and exact
integer ciphertext language. The fresh-sign byte decoder recovers the
original bits for positive scaling.
&
\lean{HE.HCP.CoupledDecoding.decodeCoefficients_realize}\par
\lean{HE.Protocol.Storage.CoupledRecovery.recoverWithKeyBytes_encode_sample}\par
\lean{HE.Protocol.Storage.CoupledRecovery.recover_eq_some_iff_sampleBytes}\par
\lean{HE.Protocol.Transfer.PathSignature.NativeSignSampler.recoverBitsWithKeyBytes_ciphertextBytes}
\\\midrule
\textbf{P1-N05. Fresh one-use family.}\newline
\Cref{sec:coupled-sampling,sec:coupled-finite-carrier}; actual canonical
ciphertext law and the original sampler's scalar recovery diagnostic.
&
\lean{HE.Protocol.Transfer.PathSignature.NativeSignSampler.ciphertextDistribution_eq}\par
\lean{HE.Protocol.Transfer.PathSignature.NativeSignPrivacy.ciphertext_guess_probability}\par
\lean{HE.Protocol.Transfer.PathSignature.NativeRecovery.uniformSeed_recovery}
\\\midrule
\textbf{P1-N06. Finite instance and private operator opening.}\newline
\Cref{prop:coupled-committed-recovery,thm:coupled-instance-recovery}.
Canonical checks retain the original coefficients or exhibit complete-preimage
collisions; the fixed-hash transfer names its actual distinguishing test.
&
\lean{HE.Protocol.Storage.OperatorCommitment.matrix_eq_or_collision}\par
\lean{HE.Protocol.Transfer.PathSignature.NativeCommittedPrivacy.sampleWitness_bounded}\par
\lean{HE.Protocol.Transfer.PathSignature.NativeCommitted.extracted_target_or_collisions}\par
\lean{HE.Protocol.Transfer.PathSignature.NativeCommittedPrivacy.recovery_le_guess_plus_operatorHiding}
\\\midrule
\textbf{P1-N07. Bounded-query recovery.}\newline
\Cref{thm:coupled-native-rom}; fresh empty-cache shared-oracle experiment,
public-only decoder and original partial integer coefficients.
&
\lean{HE.Protocol.Storage.OperatorCommitment.RandomOracle.hit_probability_le}\par
\lean{HE.Protocol.Storage.OperatorCommitment.RandomOracle.event_probability_le}\par
\lean{HE.Protocol.Transfer.PathSignature.NativeCommittedRandomOracle.game_query_bound}\par
\lean{HE.Protocol.Transfer.PathSignature.NativeCommittedRandomOracle.ideal_recovery_probability}\par
\lean{HE.Protocol.Transfer.PathSignature.NativeCommittedRandomOracle.recovery_probability_le}\par
\lean{HE.Protocol.Transfer.PathSignature.NativeCommittedRandomOracle.partial_coeff_recovery_probability_le}
\\\midrule
\textbf{P1-N08. Strict native CVP comparison.}\newline
\Cref{thm:strict-native-cvp-separation}; one-call recovery with the supplied
lattice and a strict public recovery gap with the same fixed oracle.
&
\lean{HE.HCP.CoupledRecoverySeparation.knownLattice_recovers_bits}\par
\lean{HE.HCP.CoupledRecoverySeparation.recovery_comparison}\par
\lean{HE.HCP.CoupledRecoverySeparation.strict_recovery_gap}
\\\midrule
\textbf{P1-N09. Harmonic lifts of arbitrary CVP targets.}\newline
\Cref{prop:hodge-projection-is-cvp,prop:coupled-cvp-residual,prop:unrestricted-coupled-cvp-lift,cor:coupled-cvp-np-complete}.
Lean checks the basis-span lift, the exact residual split and the search and
decision transfers under the scale inequality. The paper derives the scale
and the complexity result.
&
\lean{HarmonicCVPBridge.solution_iff_projectedInstance}\par
\lean{HarmonicCVPBridge.oracleCorrect_of_harmonicTargets}\par
\lean{HE.HCP.CoupledCVP.squaredDistance_eq_residual_add_gram}\par
\lean{HE.HCP.CoupledCVPLift.source_solution_of_scaled_solution}\par
\lean{HE.HCP.CoupledCVPLift.source_threshold_iff_scaledDistance}
\\\midrule
\textbf{P1-N10. Compatible answers at one ciphertext.}\newline
\Cref{prop:coupled-cvp-candidate-count}; exact counts for coefficients and
operative lattices, and actual CVP recovery for each matching key.
&
\lean{HE.HCP.CoupledRecoverySeparation.compatible_coefficients_card}\par
\lean{HE.HCP.CoupledRecoverySeparation.compatible_lattices_card}\par
\lean{HE.HCP.CoupledRecoverySeparation.knownLattice_recovers_masked_bits}
\\\midrule
\textbf{P1-N11. Chordal geometry.}\newline
\Cref{prop:coupled-chordal-distance}; actual orthogonal projectors and
the distances between compatible fresh-sign planes.
&
\lean{HE.Laplacian.Grassmannian.chordalDistance_eq_zero_iff}\par
\lean{HE.HCP.CoupledProjectorMetric.projectorMatrix_eq_graph}\par
\lean{HE.HCP.CoupledProjectorMetric.chordalDistance_sq_eq_hamming}
\\\midrule
\textbf{P1-N12. Minimum residual and sharp recovery gap.}\newline
\Cref{prop:fresh-sign-cvp-residual}; equal minimum squared residuals,
the exact next cost and the strict additive tolerance on the supplied lattice.
&
\lean{HE.HCP.CoupledSignResidual}\par
\lean{squaredDistance_fixed_mask}\par
\lean{cost_at_message}\par
\lean{cost_gap_attained}\par
\lean{recover_of_cost_lt}
\\\midrule
\textbf{P1-N13. Exact recovery under finite disclosure.}\newline
\Cref{prop:finite-disclosure}; sharp target-multiplicity bound at a fixed
observation and under deterministic disclosure. The canonical ciphertext
fiber, sign-coordinate disclosures and constant minimum cost specialize
the bound to the fresh-sign family.
&
\lean{HE.HCPBridge.fixedObservation_optimal_probability}\par
\lean{HE.HCP.FiniteRecovery}\par
\lean{optimalProbability}, \lean{successCount_le_optimalCount},\par
\lean{exists_successCount_eq_optimalCount}\par
\lean{HE.HCP.FiniteRecovery.Native}\par
\lean{fixedMaskFiberEquiv}, \lean{uniformKey_success_eq},\par
\lean{optimalProbability_eq_view_count},\par
\lean{optimalProbability_le_bit_budget},\par
\lean{optimalProbability_eq_coordinate_disclosure},\par
\lean{optimalProbability_add_minimum_cost}
\\\midrule
\textbf{P1-N14. Bounded-query recovery with key disclosure.}\newline
\Cref{cor:coupled-native-rom-disclosure}; deterministic disclosure of the
sign key in the same one-use byte-oracle experiment. The ideal guessing
bound uses the number of disclosure values and retains the original nonce
term.
&
\lean{HE.Protocol.Transfer.PathSignature.NativeCommittedRandomOracle}\par
\lean{uniform_guess_with_disclosure_range_le}\par
\lean{Disclosure.ideal_game_masked}\par
\lean{Disclosure.ideal_recovery_probability_le}\par
\lean{Disclosure.recovery_probability_le}\par
\lean{Disclosure.partial_coeff_recovery_probability_le}\par
\lean{Disclosure.bit_disclosure_recovery_probability_le}
\\\midrule
\textbf{P1-N15. Bounded ciphertext circuit.}\newline
\Cref{prop:coupled-field-refinement}. Honest range assignments and a unique
bounded lift of the core wires to the original integer equations.
&
\lean{HE.Protocol.Transfer.PathSignature}\par
\lean{NativeCircuitRanges}\par
\lean{nativeWitness_checked}, \lean{existsUnique_integer_assignment},\par
\lean{exists_integer_observed}
\\\midrule
\textbf{P1-N16. Fresh-sign decryption and reuse.}\newline
\Cref{prop:coupled-sign-checked-recovery,prop:coupled-sign-reuse}.
Exact accepted byte language, messages under alternative valid keys,
reused-key agreement and compatible-pair count, and known-message recovery.
&
\lean{HE.Protocol.Transfer.PathSignature.NativeSignSampler}\par
\lean{CheckedRecovery.recover_eq_some_iff_ciphertextBytes}\par
\lean{Reuse.recoverBitsWithKeyBytes_alternativeSignKey}\par
\lean{Reuse.recoverAgreement_ciphertextBytes}\par
\lean{Reuse.compatible_message_pairs_card}\par
\lean{Reuse.recoverKnownKey_ciphertextBytes}\par
\lean{Reuse.recoverReusedMessage_ciphertextBytes}
\\\midrule
\textbf{P1-N17. One-use chosen-message privacy.}\newline
\Cref{thm:coupled-native-rom-challenge}. Fixed pair and adaptive postrelease
queries to the native byte oracle.
&
\lean{HE.Protocol.Transfer.PathSignature.NativeCommittedRandomOracle}\par
\lean{Indistinguishability.ideal_guess_probability}\par
\lean{Indistinguishability.guess_probability_le}\par
\lean{Indistinguishability.half_le_guess_probability_add}
\\
\end{longtable}
\endgroup

\subsection{Historical Palomar comparison}\label{app:historical-palomar}

The Palomar submission used an earlier revision of the manuscript.
Its archived comparator configuration selects 16 manuscript-facing and 20
supporting declarations across six Hodge/CVP groups from that revision.
\Cref{tab:paper1-parity,tab:paper1-support,tab:paper1-definitions} retain the
submission's crosswalk, including canonical consequences and the premises of
the selected statements. The source pin and verification record in
\lean{lean/Palomar/PaperOne/README.md} apply to that historical selection.
The recorded axiom audit for those declarations allows \lean{propext},
\lean{Quot.sound} and \lean{Classical.choice}.

In these tables, \lean{E} abbreviates \lean{HE.Palomar.PaperOne},
\lean{CS} abbreviates \lean{E.CellularSheaf}, \lean{HB} abbreviates
\lean{E.HCPBridge}, and \lean{AF} abbreviates \lean{E.AxisFamily}.

\begingroup
\small
\setlength{\tabcolsep}{6pt}
\renewcommand{\arraystretch}{1.12}
\renewcommand{\lean}[1]{\begingroup\footnotesize\Urlmuskip=0mu plus 1mu\path{#1}\endgroup}
\begin{longtable}{@{}>{\raggedright\arraybackslash}p{0.41\linewidth}>{\raggedright\arraybackslash}p{0.55\linewidth}@{}}
\caption{Focused Hodge/CVP statement groups}\label{tab:paper1-parity}\\
\toprule
\textbf{Statement and scope} & \textbf{Formal declarations or boundary}\\
\midrule
\endfirsthead
\caption[]{Focused Hodge/CVP statement groups (continued).}\\
\toprule
\textbf{Statement and scope} & \textbf{Formal declarations or boundary}\\
\midrule
\endhead
\midrule
\multicolumn{2}{r@{}}{\emph{Continued on the next page.}}\\
\endfoot
\bottomrule
\endlastfoot

\textbf{P1-S01.}\newline
\Cref{prop:hodge-decomposition}; private recovery in Sections~\ref{sec:primitives}--\ref{sec:cvp-comparison}.
The complex equation and actual adjoint Hodge operator give componentwise
unique harmonic, exact and coexact parts. Integer coefficient recovery
is classical and noncomputable in the supplied integral harmonic lattice.
&
\lean{E.Hodge_Decomposition_Theorem}\par
\lean{E.Hodge_Decomposition_unique}\par
\lean{CS.HodgeEncryptedCochain.harmonicOneProjection_observed}\par
\lean{CS.HarmonicLatticePresentation.decodeCoefficients_encrypt}
\\\midrule

\textbf{P1-S02.}\newline
\Cref{prop:hcp-to-cvp}.
A field-linear image of the displayed generator is realized as a kernel.
The consequence concerns that image, not integer-lattice membership.
&
\lean{E.exists_matrix_ker_eq_range}\par
\lean{E.Hodge_GeneratedSubspace_Reduction_Proof}
\\\midrule

\textbf{P1-S03.}\newline
\Cref{prop:hodge-projection-is-cvp}.
The discrete lattice is explicitly supplied and lies in the projected
subspace. The Pythagorean identity preserves exactly its closest-vector answers.
&
\lean{E.squaredDistance_projection_decomposition}\par
\lean{E.solution_iff_projectedTarget}
\\\midrule

\textbf{P1-S04.}\newline
Cellular specialization of
\cref{prop:hodge-projection-is-cvp}.
For the actual Hodge sample and supplied harmonic message lattice, the
operative message is the unique closest lattice point.
&
\lean{E.concreteSheafHCP_CVP_solution_iff_message}\par
\lean{E.concreteSheafHCP_CVP_existsUnique}
\\\midrule

\textbf{P1-S05.}\newline
\Cref{prop:cvp-insufficient,prop:cellular-axis-ambiguity,cor:lattice-member-separation}.
Generic separation retains closed-test invariance and kernel separation;
oracle corollaries also retain existence of a correct explicit-lattice
CVP oracle. The constructed integral family has one public observation
and distinct nonzero targets. Its supplied rank-one lattice is solvable by
that same correct oracle; shared-class recovery is impossible, as is exact
recovery in every compatible family retaining the separated pair.
&
\lean{HB.operativeTargets_ne_of_kernelSeparation}\par
\lean{HB.cvpOracle_insufficient_of_kernelSeparation}\par
\lean{HB.cvpOracle_insufficient_for_operativeTarget}\par
\lean{AF.cellular_axis_ambiguity}\par
\lean{AF.cvp_oracle_insufficient}\par
Canonical consequences:\par
\lean{HE.HCPBridge.no_solver_exactOnOperativeTarget_on_family}\par
\lean{HE.HCP.AxisFamily.cvp_known_member_hidden_member_separation}
\\\midrule

\textbf{P1-S06.}\newline
Weighted example in \cref{sec:axis-family}.
The common observation is \(b=(1,2)\), with private integer coefficients
two and one. Each privately presented lattice has its own unique CVP answer.
&
\lean{AF.weighted_axis_recovery_obstruction}
\\
\end{longtable}

\begin{longtable}{@{}>{\raggedright\arraybackslash}p{0.41\linewidth}>{\raggedright\arraybackslash}p{0.55\linewidth}@{}}
\caption{Supporting statements}\label{tab:paper1-support}\\
\toprule
\textbf{Statement and scope} & \textbf{Formal declarations or boundary}\\
\midrule
\endfirsthead
\caption[]{Supporting statements (continued).}\\
\toprule
\textbf{Statement and scope} & \textbf{Formal declarations or boundary}\\
\midrule
\endhead
\midrule
\multicolumn{2}{r@{}}{\emph{Continued on the next page.}}\\
\endfoot
\bottomrule
\endlastfoot

\textbf{Complexes and semantics.}\newline
Sections~\ref{sec:primitives}--\ref{sec:cvp-comparison}.
Actual coboundary composition, the constructed complex isomorphism and
the corresponding three-term differential.
&
\lean{CS.delta1_comp_delta0}\par
\lean{AF.cellular_complex_iso}\par
\lean{E.ThreeTermData.differential_comp}\par
\lean{AF.boundarySqZero}
\\\midrule

\textbf{Axis coordinates.}\newline
\cref{sec:axis-family}.
Coordinate vectors and the distinct complementary-axis choices in the examples.
&
\lean{AF.unitCochain_apply}\par
\lean{AF.other_zero}\par
\lean{AF.other_one}\par
\lean{AF.other_ne}
\\\midrule

\textbf{Axis Hodge operators.}\newline
\cref{sec:axis-family}.
Both coboundaries, the adjoint, the degree-one Laplacian and its harmonicity criterion.
&
\lean{AF.delta0_eq_zero}\par
\lean{AF.delta1_apply}\par
\lean{AF.delta1_adjoint_eq}\par
\lean{AF.degreeOneLaplacian_apply}\par
\lean{AF.harmonic_iff}
\\\midrule

\textbf{Integral lattice.}\newline
\cref{sec:axis-family}.
Restriction-compatible stalk integrality, harmonic basis vectors,
encoding injectivity and membership in the chosen lattice.
&
\lean{AF.standardStalkLattice_contains_iff}\par
\lean{AF.restriction_preserves}\par
\lean{AF.unitCochain_integral}\par
\lean{AF.other_unit_harmonic}\par
\lean{AF.harmonicBasis_injective}\par
\lean{AF.harmonicLatticeVector_eq}\par
\lean{AF.latticeVector_mem}
\\
\end{longtable}

\begin{longtable}{@{}>{\raggedright\arraybackslash}p{0.41\linewidth}>{\raggedright\arraybackslash}p{0.55\linewidth}@{}}
\caption{Definitions and proof boundaries}\label{tab:paper1-definitions}\\
\toprule
\textbf{Statement and scope} & \textbf{Formal declarations or boundary}\\
\midrule
\endfirsthead
\caption[]{Definitions and proof boundaries (continued).}\\
\toprule
\textbf{Statement and scope} & \textbf{Formal declarations or boundary}\\
\midrule
\endhead
\midrule
\multicolumn{2}{r@{}}{\emph{Continued on the next page.}}\\
\endfoot
\bottomrule
\endlastfoot

\textbf{HCP input and target.}\newline
Section~2. The public pair contains a classifier value and observed cochain.
A private representative selects the operative kernel and hence the
orthogonal-projection target.
&
Canonical definitions:\par
\lean{HE.HCPBridge.Observation}\par
\lean{HE.HCPBridge.HiddenHodgeWitness}\par
\lean{HE.HCP.operativeHarmonicComponent}\par
\lean{HE.HCPBridge.operativeTarget}
\\\midrule

\textbf{Explicit-lattice CVP.}\newline
\cref{sec:known-lattice-cvp}. An injective real basis presents the lattice through integer
coefficients. An instance adds a target; a solution is a closest discrete
lattice point.
&
Canonical definitions:\par
\lean{CVP.ExplicitLattice}\par
\lean{CVP.Instance}\par
\lean{CVP.Solution}
\\\midrule

\textbf{Semantic realization.}\newline
Sections~1--2. The public classifier factors through the specified closed
theory. The operative Hodge operator is separate presentation-level data;
equal classifier values need not identify it. Completeness adds reflection
of the declared homotopy equivalence; with invariance, public fibers are
exactly its classes. The constructed family proves this condition.
&
Canonical definitions and classifier theorems:\par
\lean{HE.Axioms.PublicClassifier}\par
\lean{HE.Helfer.ClosedTheory}\par
\lean{HE.HCPBridge.Realization}\par
\lean{HE.HCPBridge.ReflectsHomotopyEquivalence}\par
\lean{HE.HCPBridge.samePublicClass_iff_homotopyEquivalent}\par
\lean{HE.HCPBridge.publicClassOf_eq_setOf_homotopyEquivalent}\par
\lean{HE.HCP.AxisFamily.theory_eq_iff_homotopyEquivalent}
\\\midrule

\textbf{Hodge sample and lattice.}\newline
Sections~\ref{sec:primitives}--\ref{sec:cvp-comparison}. The sample combines an integral harmonic message with
exact/coexact noise. A supplied harmonic-lattice presentation has full real
span; it need not contain every intrinsic integral harmonic cochain.
&
Canonical definitions:\par
\lean{HE.Sheaf.CellularSheaf.HodgeEncryptedCochain}\par
\lean{HE.Sheaf.CellularSheaf.HarmonicLatticePresentation}
\\\midrule

\textbf{External theorem and semantic premise.}\newline
Sections~\ref{sec:mathematical-framework} and~\ref{sec:axis-family}; distinct from the proved homological fragment.
&
Helfer's general fibrational/path-invariance theorem is cited background.
The generic bridge uses the declared closed-test invariance; the concrete
family proves it for its homological tests. Full first-order invariance is
not a theorem of this entry.
\\\midrule

\textbf{Assumption: computational hardness.}\newline
\cref{sec:known-lattice-cvp}; distinct from the proved exact comparisons.
&
Known-lattice CVP hardness is an external computational interface.
The selected algebraic comparisons and non-identifiability examples retain
their exact-recovery scope. The coupled native random-oracle theorem above
is a separate sampled-query result.
\\
\end{longtable}

\endgroup
